\ifdefined\gkver\else\def\gkver{student}\fi
\documentclass[\gkver]{gaokaozhenti}
\begin{document}

\chapter{2009年安徽理综}
%% paperType: 理综物理部分

\begin{enumerate}
\item
%% number: 14
%% paperName: 2009年安徽理综第 14 题 6 分
%% typeId: 1
%% type: 单选题
%% chapter: 原子和原子核
%% point: 核聚变
%% method: 无
%% score: 6
%% degree: 650
%% duplicateId: 0
%% body:
原子核聚变可望给人类未来提供丰富的洁净能源。当氘等离子体被加热到适当高温时，氘核参与的几种聚变反应可能发生，放出能量。这几种反应总的效果可以表示为

\[6\ce{^{2}_{1}H}\rightarrow k\ce{^{4}_{2}He}+d\ce{^{1}_{1}H}+2\ce{^{1}_{0}n}+43.15\UMeV\]

由平衡条件可知\xzanswer{B}

\fourchoices[answer=B]
{$k=1$，$d=4$}
{$k=2$，$d=2$}
{$k=1$，$d=6$}
{$k=2$，$d=3$}

%% answer: B
%% memo:
\memoanswer{
由质量数守恒和电荷数守恒，分别有 $4k+d=10$，$2k+d=6$，解得 $k=2$，$d=2$。正确选项为 B。
}

\item
%% number: 15
%% paperName: 2009年安徽理综第 15 题 6 分
%% typeId: 1
%% type: 单选题
%% chapter: 曲线运动
%% point: 人造地球卫星
%% method: 无
%% score: 6
%% degree: 650
%% duplicateId: 0
%% body:
2009 年 2 月 11 日，俄罗斯的“宇宙-2251”卫星和美国的“铱-33”卫星在西伯利亚上空约 $805\Ukm$ 处发生碰撞。这是历史上首次发生的完整在轨卫星碰撞事件。碰撞过程中产生的大量碎片可能会影响太空环境。假定有甲、乙两块碎片，绕地球运动的轨道都是圆，甲的运行速率比乙的大，则下列说法中正确的是\xzanswer{D}

\fourchoices[answer=D]
{甲的运行周期一定比乙的长}
{甲距地面的高度一定比乙的高}
{甲的向心力一定比乙的小}
{甲的加速度一定比乙的大}

%% answer: D
%% memo:
\memoanswer{
由 $v=\sqrt{\frac{GM}{r}}$ 可知，甲的速率大，甲碎片的轨道半径小，故 B 错；由公式 $T=2\pi\sqrt{\frac{R^{3}}{GM}}$ 可知甲的周期小，故 A 错；由于未知两碎片的质量，无法判断向心力的大小，故 C 错；碎片的加速度是指引力加速度，由 $\frac{GMm}{R^{2}}=ma$ 得 $\frac{GM}{R^{2}}=a$，可知甲的加速度比乙大，故 D 对。
}

\item
%% number: 16
%% paperName: 2009年安徽理综第 16 题 6 分
%% typeId: 1
%% type: 单选题
%% chapter: 直线运动
%% point: 运动图象
%% method: 无
%% score: 6
%% degree: 450
%% duplicateId: 0
%% body:
大爆炸理论认为，我们的宇宙起源于 137 亿年前的一次大爆炸。除开始瞬间外，在演化至今的大部分时间内，宇宙基本上是匀速膨胀的。上世纪末，对 1A 型超新星的观测显示，宇宙正在加速膨胀，面对这个出人意料的发现，宇宙学家探究其背后的原因，提出宇宙的大部分可能由暗能量组成，它们的排斥作用导致宇宙在近段天文时期内开始加速膨胀。如果真是这样，则标志宇宙大小的宇宙半径 $R$ 和宇宙年龄 $t$ 的关系，大致是下面哪个图像？\xzanswer{C}

\fourchoices[answer=C, ispicture=true, h=2.6cm]
{figs/16a.png}
{figs/16b.png}
{figs/16c.png}
{figs/16d.png}

%% answer: C
%% memo:
\memoanswer{
图像中的纵坐标宇宙半径 $R$ 可以看作是星球发生的位移 $x$，因而其切线的斜率就是宇宙半径增加的快慢程度。由题意，宇宙加速膨胀，其半径增加的速度越来越大。故选 C。
}

\item
%% number: 17
%% paperName: 2009年安徽理综第 17 题 6 分
%% typeId: 1
%% type: 单选题
%% chapter: 运动和力
%% point: 牛顿第二定律
%% method: 无
%% score: 6
%% degree: 450
%% duplicateId: 0
%% body:
为了节省能量，某商场安装了智能化的电动扶梯。无人乘行时，扶梯运转得很慢；有人站上扶梯时，它会先慢慢加速，再匀速运转。一顾客乘扶梯上楼，恰好经历了这两个过程，如图所示。那么下列说法中正确的是\xzanswer{C}

\onepicture[w=6.0cm]{figs/17.png}

\fourchoices[answer=C]
{顾客始终受到三个力的作用}
{顾客始终处于超重状态}
{顾客对扶梯作用力的方向先指向左下方，再竖直向下}
{顾客对扶梯作用的方向先指向右下方，再竖直向下}

%% answer: C
%% memo:
\memoanswer{
在慢慢加速的过程中顾客受到的摩擦力水平向左，电梯对其的支持力和摩擦力的合力方向指向右上，由牛顿第三定律，它的反作用力即人对电梯的作用方向指向向左下；在匀速运动的过程中，顾客与电梯间的摩擦力等于零，顾客对扶梯的作用仅剩下压力，方向沿竖直向下。

\onepicture[w=4.5cm]{figs/17_an.jpg}
}

\item
%% number: 18
%% paperName: 2009年安徽理综第 18 题 6 分
%% typeId: 1
%% type: 单选题
%% chapter: 电场
%% point: 带电粒子的往复运动
%% method: 无
%% score: 6
%% degree: 450
%% duplicateId: 0
%% body:
在光滑的绝缘水平面上，有一个正方形的 abcd，顶点 a、c 处分别固定一个正点电荷，电荷量相等，如图所示。若将一个带负电的粒子置于 b 点，自由释放，粒子将沿着对角线 bd 往复运动。粒子从 b 点运动到 d 点的过程中\xzanswer{D}

\onepicture[w=4.0cm]{figs/18.png}

\fourchoices[answer=D]
{先作匀加速运动，后作匀减速运动}
{先从高电势到低电势，后从低电势到高电势}
{电势能与机械能之和先增大，后减小}
{电势能先减小，后增大}

%% answer: D
%% memo:
\memoanswer{
由于负电荷受到的电场力是变力，加速度是变化的。所以 A 错；

由等量正电荷的电场分布知道，在两电荷 ac 连线的中垂线 bd 的交点 O 点的电势最高，所以从 b 到 a，电势是先增大后减小，故 B 错；

由于只有电场力做功，所以只有电势能与动能的相互转化，故电势能与机械能之和守恒，C 错；

由 b 到 O 电场力做正功，电势能减小，由 O 到 d 电场力做负功，电势能增加，D 对。

\onepicture[w=4.2cm]{figs/18_an.jpg}
}

\item
%% number: 19
%% paperName: 2009年安徽理综第 19 题 6 分
%% typeId: 1
%% type: 单选题
%% chapter: 磁场
%% point: 带电粒子在磁场中的运动
%% method: 无
%% score: 6
%% degree: 450
%% duplicateId: 0
%% body:
右图是科学史上一张著名的实验照片，显示一个带电粒子在云室中穿过某种金属板运动的径迹。云室旋转在匀强磁场中，磁场方向垂直照片向里。云室中横放的金属板对粒子的运动起阻碍作用。分析此径迹可知粒子\xzanswer{A}

\onepicture[w=5.5cm]{figs/19.jpg}

\fourchoices[answer=A]
{带正电，由下往上运动}
{带正电，由上往下运动}
{带负电，由上往下运动}
{带负电，由下往上运动}

%% answer: A
%% memo:
\memoanswer{
粒子穿过金属板后，速度变小，由半径公式 $r=\frac{mv}{qB}$ 可知，半径变小，粒子运动方向为由下向上；又由于洛伦兹力的方向指向圆心，由左手定则，粒子带正电。选 A。
}

\item
%% number: 20
%% paperName: 2009年安徽理综第 20 题 6 分
%% typeId: 1
%% type: 单选题
%% chapter: 电磁感应
%% point: 电磁感应的电量
%% method: 无
%% score: 6
%% degree: 450
%% duplicateId: 0
%% body:
如图甲所示，一个电阻为 $R$，面积为 $S$ 的矩形导线框 abcd，水平旋转在匀强磁场中，磁场的磁感应强度为 $B$，方向与 $ad$ 边垂直并与线框平面成 $45^\circ$ 角，o、o$'$ 分别是 ab 和 cd 边的中点。现将线框右半边 obco$'$ 绕 oo$'$ 逆时针 $90^\circ$ 到图乙所示位置。在这一过程中，导线中通过的电荷量是\xzanswer{A}

\onepicture[w=9.5cm]{figs/20.png}

\fourchoices[answer=A]
{$\frac{\sqrt{2}BS}{2R}$}
{$\frac{\sqrt{2}BS}{R}$}
{$\frac{BS}{R}$}
{0}

%% answer: A
%% memo:
\memoanswer{
对线框的右半边（obco$'$）未旋转时整个回路的磁通量 $\Phi_{1}=BS\sin 45^{\circ}=\frac{\sqrt{2}}{2}BS$。

对线框的右半边（obco$'$）旋转 $90^{\circ}$ 后，穿进跟穿出的磁通量相等，如右图整个回路的磁通量 $\Phi_{2}=0$。$\Delta\Phi=\Phi_{2}-\Phi_{1}=\frac{\sqrt{2}}{2}BS$。根据公式 $q=\frac{\Delta\Phi}{R}=\frac{\sqrt{2}BS}{2R}$。选 A。

\onepicture[w=4.0cm]{figs/20_an1.jpg}
\onepicture[w=4.0cm]{figs/20_an2.jpg}
}

\item
%% number: 21
%% paperName: 2009年安徽理综第 21 题 6 分
%% typeId: 4
%% type: 实验题
%% chapter: 功和能
%% point: 验证动能定理
%% method: 无
%% score: 6
%% degree: 450
%% duplicateId: 0
%% body:
\begin{enumerate}
	\item
	用多用电表进行了几次测量，指针分别处于 a、b 的位置，如图所示。若多用电表的选择开关处于下面表格中所指的档位，a 和 b 的相应读数是多少？请填在表格中。

	\onepicture[w=8.0cm]{figs/21a.png}

	\begin{center}
	\begin{tabular}{|c|c|c|}
	\hline
	指针位置 & 选择开关所处挡位 & 读数 \\ \hline
	a & 直流电流 $100\UmA$ & \tkanswer[1.5cm]{}$\UmA$ \\ \hline
	a & 直流电压 $2.5\UV$ & \tkanswer[1.5cm]{}$\UV$ \\ \hline
	b & 电阻 $\times100$ & \tkanswer[1.5cm]{}$\UO$ \\ \hline
	\end{tabular}
	\end{center}

	\item
	用右图所示的电路，测定一节干电池的电动势和内阻。电池的内阻较小，为了防止在调节滑动变阻器时造成短路，电路中用一个定值电阻 $R_{0}$ 起保护作用。除电池、开关和导线外，可供使用的实验器材还有：

	（a）电流表（量程 $0.6\UA$、$3\UA$）；

	（b）电压表（量程 $3\UV$、$15\UV$）；

	（c）定值电阻（阻值 $1\UO$、额定功率 $5\UW$）；

	（d）定值电阻（阻值 $10\UO$，额定功率 $10\UW$）；

	（e）滑动变阻器（阻值范围 $0\sim10\UO$、额定电流 $2\UA$）；

	（f）滑动变阻器（阻值范围 $0\sim100\UO$、额定电流 $1\UA$）。

	那么

	（1）要正确完成实验，电压表的量程应选择\tkanswer[1.2cm]{}$\UV$，电流表的量程应选择\tkanswer[1.2cm]{}$\UA$；$R_{0}$ 应选择\tkanswer[1.2cm]{}$\UO$ 的定值电阻，$R$ 应选择阻值范围是\tkanswer[1.2cm]{}$\UO$ 的滑动变阻器。

	（2）引起该实验系统误差的主要原因是\tkanswer[6cm]{}。

	\onepicture[w=4.5cm, align=right]{figs/21b.png}

	\item
	探究力对原来静止的物体做的功与物体获得的速度的关系，实验装置如图所示，实验主要过程如下：

	\begin{steps}
		\item
		设法让橡皮筋对小车做的功分别为 $W$、$2W$、$3W$、……；
		\item
		分析打点计时器打出的纸带，求出小车的速度 $v_{1}$、$v_{2}$、$v_{3}$、……；
		\item
		作出 $W-v$ 草图；
		\item
		分析 $W-v$ 图像。如果 $W-v$ 图像是一条直线，表明 $W\propto v$；如果不是直线，可考虑是否存在 $W\propto v^{2}$、$W\propto v^{3}$、$W\propto\sqrt{v}$ 等关系。
	\end{steps}

	以下关于该试验的说法中有一项不正确，它是\tkanswer[1.2cm]{}。

	（A）本实验设法让橡皮筋对小车做的功分别为 $W$、$2W$、$3W$、……。所采用的方法是选用同样的橡皮筋，并在每次实验中使橡皮筋拉伸的长度保持一致。当用 1 条橡皮筋进行实验时，橡皮筋对小车做的功为 $W$，用 2 条、3 条、……橡皮筋并在一起进行第 2 次、第 3 次、……实验时，橡皮筋对小车做的功分别是 $2W$、$3W$、……。

	（B）小车运动中会受到阻力，补偿的方法，可以使木板适当倾斜。

	（C）某同学在一次实验中，得到一条记录纸带。纸带上打出的点，两端密、中间疏。出现这种情况的原因，可能是没有使木板倾斜或倾角太小。

	（D）根据记录纸带上打出的点，求小车获得的速度的方法，是以纸带上第一点到最后一点的距离来进行计算。

	\onepicture[w=8.0cm, align=right]{figs/21c.png}
\end{enumerate}

%% answer:
\jdanswer{
\begin{enumerate}
	\item
	$23\UmA$，$0.57\UV$，$320\UO$。
	\item
	（1）$3\UV$，$0.6\UA$，$1\UO$，$0\sim10\UO$；（2）由于电压表的分流作用造成电流表读数总是比电池实际输出电流小。
	\item
	D。
\end{enumerate}
}

%% memo:
\memoanswer{
直流电流 $100\UmA$ 档读第二行“$0\sim10$”一排，最小度值为 $2\UmA$，估读到 $1\UmA$ 就可以了；直流电压 $2.5\UV$ 档读第二行“$0\sim250$”一排，最小分度值为 $0.05\UV$，估读到 $0.01\UV$ 就可以了；电阻 $\times100$ 档读第一行，测量值等于表盘上读数“3.2”乘以倍率“100”。
}

\memoanswer{
由于电源是一节干电池（$1.5\UV$），所选量程为 $3\UV$ 的电压表；估算电流时，考虑到干电池的内阻一般几 $\UO$ 左右，加上保护电阻，最大电流在 $0.5\UA$ 左右，所以选量程为 $0.6\UA$ 的电流表；由于电池内阻很小，所以保护电阻不宜太大，否则会使得电流表、电压表取值范围小，造成的误差大；滑动变阻器的最大阻值一般比电池内阻大几倍就好了，取 $0\sim10\UO$ 能很好地控制电路中的电流和电压，若取 $0\sim100\UO$ 会出现开始几乎不变最后突然变化的现象。

关于系统误差一般由测量工具和所造成测量方法造成的，一般具有倾向性，总是偏大或者偏小。本实验中由于电压表的分流作用造成电流表读数总是比测量值小，造成 $E_{\text{测}}<E_{\text{真}}$，$r_{\text{测}}<r_{\text{真}}$。
}

\memoanswer{
本实验的目的是探究橡皮绳做的功与物体获得速度的关系。这个速度是指橡皮绳做功完毕时的速度，而不整个过程的平均速度，所以 D 选项是错误的。
}

\item
%% number: 22
%% paperName: 2009年安徽理综第 22 题 14 分
%% typeId: 6
%% type: 计算题
%% chapter: 运动和力
%% point: 连接体
%% method: 无
%% score: 14
%% degree: 650
%% duplicateId: 0
%% body:
在 2008 年北京残奥会开幕式上，运动员手拉绳索向上攀登，最终点燃了主火炬，体现了残疾运动员坚忍不拔的意志和自强不息的精神。为了探究上升过程中运动员与绳索和吊椅间的作用，可将过程简化。一根不可伸缩的轻绳跨过轻质的定滑轮，一端挂一吊椅，另一端被坐在吊椅上的运动员拉住，如图所示。设运动员的质量为 $65\Ukg$，吊椅的质量为 $15\Ukg$，不计定滑轮与绳子间的摩擦。重力加速度取 $g=10\Umsq$。当运动员与吊椅一起正以加速度 $a=1\Umsq$ 上升时，试求：

\begin{enumerate}
	\item
	运动员竖直向下拉绳的力；
	\item
	运动员对吊椅的压力。
\end{enumerate}
\onepicture[w=4.0cm, align=right]{figs/22.png}

%% answer:
\jdanswer{
\begin{enumerate}
	\item
	$440\UN$
	\item
	$275\UN$
\end{enumerate}
}

%% memo:
\memoanswer{
解法一：（1）设运动员受到绳向上的拉力为 $F$，由于跨过定滑轮的两段绳子拉力相等，吊椅受到绳的拉力也是 $F$。对运动员和吊椅整体进行受力分析如图所示，则有：

$2F-(m_{\text{人}}+m_{\text{椅}})g=(m_{\text{人}}+m_{\text{椅}})a$

$F=440\UN$。

由牛顿第三定律，运动员竖直向下拉绳的力 $F'=440\UN$。

\onepicture[w=4.2cm]{figs/22_an1.jpg}

（2）设吊椅对运动员的支持力为 $F_{N}$，对运动员进行受力分析如图所示，则有：

$F+F_{N}-m_{\text{人}}g=m_{\text{人}}a$

$F_{N}=275\UN$。

由牛顿第三定律，运动员对吊椅的压力也为 $275\UN$。

\onepicture[w=4.2cm]{figs/22_an2.jpg}

解法二：设运动员和吊椅的质量分别为 $M$ 和 $m$；运动员竖直向下的拉力为 $F$，对吊椅的压力大小为 $F_{N}$。根据牛顿第三定律，绳对运动员的拉力大小为 $F$，吊椅对运动员的支持力为 $F_{N}$。分别以运动员和吊椅为研究对象，根据牛顿第二定律

$F+F_{N}-Mg=Ma$ ①

$F-F_{N}-mg=ma$ ②

由①②得 $F=440\UN$，$F_{N}=275\UN$。
}

\item
%% number: 23
%% paperName: 2009年安徽理综第 23 题 16 分
%% typeId: 6
%% type: 计算题
%% chapter: 电场
%% point: 带电粒子的偏转
%% method: 无
%% score: 16
%% degree: 450
%% duplicateId: 0
%% body:
如图所示，匀强电场方向沿 $x$ 轴的正方向，场强为 $E$。在 A（$d$，0）点有一个静止的中性微粒，由于内部作用，某一时刻突然分裂成两个质量均为 $m$ 的带电微粒，其中电荷量为 $q$ 的微粒 1 沿 $y$ 轴负方向运动，经过一段时间到达（0，$-d$）点。不计重力和分裂后两微粒间的作用。试求：

\begin{enumerate}
	\item
	分裂时两个微粒各自的速度；
	\item
	当微粒 1 到达（0，$-d$）点时，电场力对微粒 1 做功的瞬间功率；
	\item
	当微粒 1 到达（0，$-d$）点时，两微粒间的距离。
\end{enumerate}
\onepicture[w=5.5cm, align=right]{figs/23.png}

%% answer:
\jdanswer{
\begin{enumerate}
	\item
	$v_{1}=-\sqrt{\frac{-qEd}{2m}}$，$v_{2}=\sqrt{\frac{qEd}{2m}}$，方向沿 $y$ 正方向
	\item
	$P=qE\sqrt{\frac{-2qEd}{m}}$
	\item
	$2\sqrt{2}d$
\end{enumerate}
}

%% memo:
\memoanswer{
（1）微粒 1 在 $y$ 方向不受力，做匀速直线运动；在 $x$ 方向由于受恒定的电场力，做匀加速直线运动。所以微粒 1 做的是类平抛运动。设微粒 1 分裂时的速度为 $v_{1}$，微粒 2 的速度为 $v_{2}$，则有：

在 $y$ 方向上有 $-d=v_{1}t$；

在 $x$ 方向上有 $a=\frac{qE}{m}$，$-d=\frac{1}{2}at^{2}$，$v_{1}=-\sqrt{\frac{-qEd}{2m}}$。根号外的负号表示沿 $y$ 轴的负方向。

中性微粒分裂成两微粒时，遵守动量守恒定律，有 $mv_{1}+mv_{2}=0$，$v_{2}=-v_{1}=\sqrt{\frac{qEd}{2m}}$，方向沿 $y$ 正方向。

（2）设微粒 1 到达（0，$-d$）点时的速度为 $v$，则电场力做功的瞬时功率为 $P=qEv_{B}\cos\theta=qEv_{Bx}$，其中由运动学公式 $v_{Bx}=-\sqrt{-2ad}=-\sqrt{\frac{-2qEd}{m}}$，所以 $P=qE\sqrt{\frac{-2qEd}{m}}$。

（3）两微粒的运动具有对称性，如图所示，当微粒 1 到达（0，$-d$）点时发生的位移 $S_{1}=\sqrt{2}d$，则当微粒 1 到达（0，$-d$）点时，两微粒间的距离为 $BC=2S_{1}=2\sqrt{2}d$。

\onepicture[w=5.0cm]{figs/23_an.jpg}
}

\item
%% number: 24
%% paperName: 2009年安徽理综第 24 题 20 分
%% typeId: 6
%% type: 计算题
%% chapter: 曲线运动
%% point: 竖直平面圆周运动
%% method: 分情况讨论
%% score: 20
%% degree: 250
%% duplicateId: 0
%% body:
过山车是游乐场中常见的设施。下图是一种过山车的简易模型，它由水平轨道和在竖直平面内的三个圆形轨道组成，B、C、D 分别是三个圆形轨道的最低点，B、C 间距与 C、D 间距相等，半径 $R_{1}=2.0\Um$、$R_{2}=1.4\Um$。一个质量为 $m=1.0\Ukg$ 的小球（视为质点），从轨道的左侧 A 点以 $v_{0}=12.0\Ums$ 的初速度沿轨道向右运动，A、B 间距 $L_{1}=6.0\Um$。小球与水平轨道间的动摩擦因数 $\mu=0.2$，圆形轨道是光滑的。假设水平轨道足够长，圆形轨道间不相互重叠。重力加速度取 $g=10\Umsq$，计算结果保留小数点后一位数字。试求：

\begin{enumerate}
	\item
	小球在经过第一个圆形轨道的最高点时，轨道对小球作用力的大小；
	\item
	如果小球恰能通过第二圆形轨道，B、C 间距 $L$ 应是多少；
	\item
	在满足（2）的条件下，如果要使小球不能脱离轨道，在第三个圆形轨道的设计中，半径 $R_{3}$ 应满足的条件；小球最终停留点与起点 A 的距离。
\end{enumerate}
\onepicture[w=11.0cm, align=right]{figs/24.png}

%% answer:
\jdanswer{
\begin{enumerate}
	\item
	$10.0\UN$
	\item
	$12.5\Um$
	\item
	当 $0<R_{3}\leq0.4\Um$ 时，$L'=36.0\Um$；当 $1.0\Um\leq R_{3}\leq27.9\Um$ 时，$L''=26.0\Um$
\end{enumerate}
}

%% memo:
\memoanswer{
（1）设小球经过第一个圆轨道的最高点时的速度为 $v_{1}$，根据动能定理

$-\mu mgL_{1}-2mgR_{1}=\frac{1}{2}mv_{1}^{2}-\frac{1}{2}mv_{0}^{2}$ ①

小球在最高点受到重力 $mg$ 和轨道对它的作用力 $F$，根据牛顿第二定律

$F+mg=m\frac{v_{1}^{2}}{R_{1}}$ ②

由①②得 $F=10.0\UN$ ③

（2）设小球在第二个圆轨道的最高点的速度为 $v_{2}$，由题意

$mg=\frac{v_{2}^{2}}{R_{2}}$ ④

$-\mu mg\left(L_{1}+L\right)-2mgR_{2}=\frac{1}{2}mv_{2}^{2}-\frac{1}{2}mv_{0}^{2}$ ⑤

由④⑤得 $L=12.5\Um$ ⑥

（3）要保证小球不脱离轨道，可分两种情况进行讨论：

I. 轨道半径较小时，小球恰能通过第三个圆轨道，设在最高点的速度为 $v_{3}$，应满足

$mg=\frac{v_{3}^{2}}{R_{3}}$ ⑦

$-\mu mg\left(L_{1}+2L\right)-2mgR_{3}=\frac{1}{2}mv_{3}^{2}-\frac{1}{2}mv_{0}^{2}$ ⑧

由⑥⑦⑧得 $R_{3}=0.4\Um$。

II. 轨道半径较大时，小球上升的最大高度为 $R_{3}$，根据动能定理

$-\mu mg\left(L_{1}+2L\right)-2mgR_{3}=0-\frac{1}{2}mv_{0}^{2}$

解得 $R_{3}=1.0\Um$。

为了保证圆轨道不重叠，$R_{3}$ 最大值应满足

$\left(R_{2}+R_{3}\right)^{2}=L^{2}+\left(R_{3}-R_{2}\right)^{2}$

解得 $R_{3}=27.9\Um$。

综合 I、II，要使小球不脱离轨道，则第三个圆轨道的半径须满足下面的条件

$0<R_{3}\leq0.4\Um$

或 $1.0\Um\leq R_{3}\leq27.9\Um$。

当 $0<R_{3}\leq0.4\Um$ 时，小球最终停留点与起始点 A 的距离为 $L'$，则

$-\mu mgL'=0-\frac{1}{2}mv_{0}^{2}$

$L'=36.0\Um$。

当 $1.0\Um\leq R_{3}\leq27.9\Um$ 时，小球最终停留点与起始点 A 的距离为 $L''$，则

$L''=L'-2(L'-L_{1}-2L)=26.0\Um$。

\onepicture[w=6.0cm]{figs/24_an.jpg}
}

\end{enumerate}

\end{document}
